\documentclass[11pt,a4paper]{article}

\usepackage[utf8]{inputenc}
\usepackage[T1]{fontenc}
\usepackage[brazilian]{babel}
\usepackage{geometry}
\geometry{top=2.5cm, bottom=2.5cm, left=2.5cm, right=2.5cm}

\usepackage{amsmath, amssymb}
\usepackage{xcolor}
\usepackage{listings}
\usepackage{tcolorbox}
\usepackage{booktabs}
\usepackage{array}
\usepackage{graphicx}
\usepackage{tikz}
\usetikzlibrary{shapes,arrows.meta,positioning,calc,fit}
\usepackage{hyperref}

\hypersetup{
    colorlinks=true,
    linkcolor=blue!70!black,
    citecolor=blue!70!black,
    urlcolor=teal
}

% Configuracao de cores para codigo
\definecolor{codegreen}{rgb}{0.13,0.55,0.13}
\definecolor{codegray}{rgb}{0.45,0.45,0.45}
\definecolor{codepurple}{rgb}{0.58,0,0.82}
\definecolor{codeblue}{rgb}{0.0,0.2,0.7}
\definecolor{backcolour}{rgb}{0.97,0.97,0.99}

\lstdefinestyle{cstyle}{
    backgroundcolor=\color{backcolour},   
    commentstyle=\color{codegreen}\itshape,
    keywordstyle=\color{codeblue}\bfseries,
    numberstyle=\tiny\color{codegray},
    stringstyle=\color{codepurple},
    basicstyle=\ttfamily\footnotesize,
    breakatwhitespace=false,         
    breaklines=true,                 
    captionpos=b,                    
    keepspaces=true,                 
    numbers=left,                    
    numbersep=7pt,                  
    showspaces=false,                
    showstringspaces=false,
    showtabs=false,                  
    tabsize=2,
    frame=single,
    rulecolor=\color{black!20}
}
\lstset{style=cstyle}

\tcbuselibrary{skins,breakable}

\newtcolorbox{destaquebox}[1]{
    colback=blue!5!white,
    colframe=blue!65!black,
    fonttitle=\bfseries,
    title=#1,
    arc=3mm,
    boxrule=1pt,
    breakable
}

\newtcolorbox{conceitobox}[1]{
    colback=green!4!white,
    colframe=green!50!black,
    fonttitle=\bfseries,
    title=#1,
    arc=3mm,
    boxrule=1pt,
    breakable
}

\begin{document}

% -------------------------------------------------------------
% CAPA / CABECALHO FORMAL
% -------------------------------------------------------------
\begin{center}
    {\large \textbf{UNIVERSIDADE FEDERAL DE UBERLÂNDIA - UFU}}\\[0.2cm]
    {\large \textbf{FACULDADE DE COMPUTAÇÃO - FACOM}}\\[0.2cm]
    {\normalsize Disciplina: Estruturas de Dados | Professor(a): ED 2026}\\[1.5cm]

    \rule{\linewidth}{1.5pt}\\[0.4cm]
    {\LARGE \textbf{DOCUMENTAÇÃO TÉCNICA E CONCEITUAL}}\\[0.3cm]
    {\Large \textbf{Sistema Clínico VidaSaudável: Integração das Quatro Estruturas Lineares Clássicas em C}}\\[0.4cm]
    \rule{\linewidth}{1.5pt}\\[1.5cm]
\end{center}

\begin{abstract}
Este documento apresenta uma análise conceitual, arquitetural e prática da implementação do sistema de atendimento hospitalar \textbf{Clínica VidaSaudável}, desenvolvido em linguagem C. O projeto atende integralmente aos requisitos curriculares de domínio das quatro estruturas de dados lineares clássicas: \textbf{Lista Simplesmente Encadeada com Cabeça} (cadastro de pacientes), \textbf{Lista Circular} (escala rotativa dos médicos), \textbf{Fila FIFO} (sala de espera) e \textbf{Pilha LIFO} (histórico de atendimentos reversível). São detalhados os conceitos de ponteiros simples e duplos, alocação dinâmica com gerenciamento rigoroso de memória, encapsulamento de Tipos Abstratos de Dados (TAD) e a mecânica interna de cada função implementada.
\end{abstract}

\vspace{0.8cm}
\tableofcontents
\newpage

% -------------------------------------------------------------
% SEÇÃO 1: OBJETIVOS E COMPETENCIAS
% -------------------------------------------------------------
\section{Objetivo Pedagógico e Competências Demonstradas}

O objetivo central deste trabalho prático é demonstrar a capacidade de projetar, modularizar e integrar em um sistema interativo coeso as quatro estruturas de dados lineares clássicas em linguagem C.

\begin{destaquebox}{Mapeamento dos Requisitos da Atividade}
\begin{enumerate}
    \item \textbf{Domínio de Ponteiros e Memória Dinâmica}:
    \begin{itemize}
        \item Uso de ponteiros simples (\texttt{T*}) para travessia, busca e leitura de estruturas.
        \item Uso de ponteiros duplos (\texttt{T**}) para funções modificadoras onde o nó inicial ou a cabeça da coleção pode mudar na função chamadora.
        \item Alocação explícita na Heap via \texttt{malloc(sizeof(...))} e liberação completa com \texttt{free()}, evitando vazamentos de memória (\textit{memory leaks}).
    \end{itemize}
    \item \textbf{Encapsulamento e Tipos Abstratos de Dados (TAD)}:
    \begin{itemize}
        \item Separação estrita entre interface (\texttt{.h}) e implementação (\texttt{.c}).
        \item Composição de estruturas de alto nível (\texttt{Fila} e \texttt{Pilha}) aproveitando internamente as listas encadeadas sem expor seus nós para o código cliente.
    \end{itemize}
    \item \textbf{Comportamentos Lineares Distintos}:
    \begin{itemize}
        \item \textbf{FIFO (First-In, First-Out)}: Fila de espera na recepção médica.
        \item \textbf{LIFO (Last-In, First-Out)}: Pilha para permitir a operação de ``Desfazer Atendimento'' (\textit{Undo}).
        \item \textbf{Circular}: Escala contínua e cíclica de plantonistas.
        \item \textbf{Linear com Sentinela/Cabeça}: Cadastro permanente de pacientes com busca e remoção arbitrária por ID.
    \end{itemize}
\end{enumerate}
\end{destaquebox}

% -------------------------------------------------------------
% SEÇÃO 2: ARQUITETURA DA CLÍNICA VIDASAUDÁVEL
% -------------------------------------------------------------
\section{Contexto e Arquitetura da Clínica VidaSaudável}

O sistema hospitalar \textbf{Clínica VidaSaudável} gerencia o fluxo de pacientes desde o acolhimento até a consulta com o médico plantonista, registrando auditoria reversível de eventos.

\subsection{As Quatro Coleções Independentes}

\begin{table}[h!]
\centering
\small
\renewcommand{\arraystretch}{1.3}
\resizebox{\textwidth}{!}{%
\begin{tabular}{@{}p{3.2cm}p{3.3cm}p{2.8cm}p{3.2cm}c@{}}
\toprule
\textbf{Coleção / Módulo} & \textbf{Estrutura Subjacente} & \textbf{Disciplina} & \textbf{Papel no Sistema} & \textbf{Complexidade} \\ \midrule
\textbf{Pacientes Cadastrados} & Lista Encadeada c/ Cabeça & Linear com Sentinela & Cadastro permanente com ID único & Busca: $O(n)$, Rem: $O(n)$ \\
\textbf{Escala de Médicos} & Lista Circular Simples & Anel Cíclico & Plantonistas rotativos & Inserir/Avançar: $O(1)$ \\
\textbf{Sala de Espera} & Fila (TAD Fila) & FIFO (Primeiro a Chegar) & Ordem de atendimento médico & Push/Pop: $O(1)$ \\
\textbf{Histórico Atendimento} & Pilha (TAD Pilha) & LIFO (Último Atendido) & Desfazer o último atendimento & Push/Pop: $O(1)$ \\ \bottomrule
\end{tabular}%
}
\caption{Quadro Comparativo das Coleções e Estruturas de Dados}
\end{table}

\subsection{Fluxo Típico de Atendimento}

O diagrama abaixo ilustra o ciclo operacional completo entre as quatro estruturas:

\begin{center}
\begin{tikzpicture}[node distance=1.8cm and 1.5cm, auto, >=Stealth]
    % Nós
    \node [rectangle, draw=blue!70, fill=blue!10, rounded corners, text width=3.2cm, align=center, minimum height=1.2cm] (cad) 
        {\textbf{1. Cadastro}\\\small Lista c/ Cabeça\\\footnotesize(ID paciente permanente)};
        
    \node [rectangle, draw=teal!70, fill=teal!10, rounded corners, text width=3.2cm, align=center, minimum height=1.2cm, right=of cad] (fila) 
        {\textbf{2. Sala de Espera}\\\small Fila FIFO\\\footnotesize(aguardando chamada)};

    \node [rectangle, draw=orange!70, fill=orange!10, rounded corners, text width=3.2cm, align=center, minimum height=1.2cm, below=of fila] (atend) 
        {\textbf{3. Consultório}\\\small Médico da Vez\\\footnotesize(Lista Circular avança)};

    \node [rectangle, draw=purple!70, fill=purple!10, rounded corners, text width=3.2cm, align=center, minimum height=1.2cm, below=of cad] (pilha) 
        {\textbf{4. Histórico}\\\small Pilha LIFO\\\footnotesize(topo = último atendido)};

    % Arestas
    \draw [->, thick] (cad) -- node [above] {\footnotesize \texttt{Opção 9: Enfileirar}} (fila);
    \draw [->, thick] (fila) -- node [right] {\footnotesize \texttt{Opção 10: Atender}} (atend);
    \draw [->, thick] (atend) -- node [below] {\footnotesize \texttt{Registra no Histórico}} (pilha);
    \draw [->, dashed, thick, red!80!black] (pilha) -- node [left] {\footnotesize \texttt{Opção 11: Desfazer Atendimento}} (fila);
\end{tikzpicture}
\end{center}

\begin{enumerate}
    \item \textbf{Cadastro}: Paciente é inserido no cadastro permanente (\texttt{Opção 1}).
    \item \textbf{Enfileiramento}: O paciente cadastrado dá entrada na sala de espera (\texttt{Opção 9}).
    \item \textbf{Atendimento}: O próximo da fila é chamado pelo médico plantonista da vez (\texttt{Opção 10}). O paciente é desenfileirado e empilhado no histórico; o médico plantonista da escala circular avança para o próximo da lista.
    \item \textbf{Desfazer Atendimento}: Se o médico precisa cancelar ou corrigir o atendimento (\texttt{Opção 11}), o paciente no topo da pilha é desempilhado e devolvido à fila de espera.
\end{enumerate}

\newpage

% -------------------------------------------------------------
% SEÇÃO 3: DETALHAMENTO DE CADA ESTRUTURA
% -------------------------------------------------------------
\subsection{Resolução de Conflitos de Nomes e Engenharia dos Cabeçalhos (\texttt{.h})}

Em linguagem C pura, não existem \textit{namespaces} (como em C++ ou C\#). Todas as funções globais e definições de \texttt{structs}/ \texttt{typedefs} compartilham um único espaço global de nomes. Quando múltiplos arquivos de cabeçalho definem tipos genéricos iguais (ex: \texttt{No}) ou funções com nomes idênticos (ex: \texttt{criar\_lista()}), a inclusão simultânea desses arquivos em um único programa integrador como \texttt{clinica.c} dispara erros fatais de redefinição no compilador GCC e colisão de símbolos no \textit{linker}.

Para resolver esse problema e permitir que as quatro estruturas convivam harmonicamente no mesmo programa, realizamos uma refatoração arquitetural nos quatro arquivos de cabeçalho:

\begin{enumerate}
    \item \textbf{\texttt{lista\_com\_cabeca.h} (Lista para Pacientes e Pilha)}:
    \begin{itemize}
        \item O tipo da estrutura foi renomeado de \texttt{No} para \texttt{ListaNo}:
\begin{lstlisting}[language=C]
typedef struct no_lista {
  int dado;
  struct no_lista *prox;
} ListaNo;
\end{lstlisting}
        \item As funções ganharam o prefixo exclusivo \texttt{lista\_}: \texttt{lista\_criar\_lista()}, \texttt{lista\_adicionar\_inicio()}, \texttt{lista\_buscar\_valor()}, \texttt{lista\_remover\_valor()}, \texttt{lista\_tamanho\_lista()}.
    \end{itemize}

    \item \textbf{\texttt{lista\_circular.h} (Lista para Médicos e Fila)}:
    \begin{itemize}
        \item O tipo da estrutura foi renomeado para \texttt{CircularNo}:
\begin{lstlisting}[language=C]
typedef struct no_circular {
  int dado;
  struct no_circular *prox;
} CircularNo;
\end{lstlisting}
        \item Adicionou-se a função \texttt{avancar\_lista(CircularNo **L)}, fundamental para rotacionar a escala dos médicos plantonistas sem destruir nós.
    \end{itemize}

    \item \textbf{\texttt{fila.h} (Sala de Espera - FIFO)}:
    \begin{itemize}
        \item Declara explicitamente a composição interna usando a lista circular:
\begin{lstlisting}[language=C]
typedef struct {
  int tam; 
  CircularNo *L; // Composicao interna opaca
} Fila;
\end{lstlisting}
        \item Expõe apenas a interface limpa do TAD: \texttt{fila\_adicionar()}, \texttt{fila\_topo()}, \texttt{fila\_remover()}, \texttt{fila\_tamanho()}.
    \end{itemize}

    \item \textbf{\texttt{pilha.h} (Histórico Reversível - LIFO)}:
    \begin{itemize}
        \item Declara explicitamente a composição interna usando a lista com cabeça:
\begin{lstlisting}[language=C]
typedef struct {
  ListaNo *L; // Composicao interna opaca
} Pilha;
\end{lstlisting}
        \item Expõe apenas a interface limpa do TAD: \texttt{pilha\_adicionar()}, \texttt{pilha\_topo()}, \texttt{pilha\_remover()}, \texttt{pilha\_tamanho()}.
    \end{itemize}
\end{enumerate}

\subsection{Lista Encadeada Simples com Cabeça (\texttt{lista\_com\_cabeca.c/.h})}

\subsubsection{Conceito do Nó Cabeça (Sentinela)}
A lista utiliza um \textbf{nó cabeça} (ou nó sentinela) alocado na inicialização. Esse nó nunca contém dados de paciente; em vez disso:
\begin{itemize}
    \item \texttt{head->dado} armazena a \textbf{quantidade total de elementos} da lista ($O(1)$ para tamanho).
    \item \texttt{head->prox} aponta para o primeiro elemento real, ou \texttt{NULL} se a lista estiver vazia.
    \item \textbf{Vantagem técnica}: Elimina o caso especial de inserir ou remover no início. O nó anterior ao primeiro sempre existe: é o próprio nó cabeça!
\end{itemize}

\begin{center}
\begin{tikzpicture}[>=Stealth, node distance=1.2cm]
    \node [draw, fill=gray!20, minimum height=0.8cm, minimum width=1.8cm] (head) {\small Cabeça: \texttt{dado=2}};
    \node [draw, fill=blue!10, minimum height=0.8cm, minimum width=1.5cm, right=of head] (n1) {\small ID: 101};
    \node [draw, fill=blue!10, minimum height=0.8cm, minimum width=1.5cm, right=of n1] (n2) {\small ID: 205};
    \node [right=0.8cm of n2] (null) {\small \texttt{NULL}};

    \draw [->, thick] (head) -- (n1);
    \draw [->, thick] (n1) -- (n2);
    \draw [->, thick] (n2) -- (null);
\end{tikzpicture}
\end{center}

\subsubsection{Implementação das Funções Chave}

\begin{lstlisting}[language=C, caption={{Criação e Inserção no Fim com Nó Cabeça}}]
ListaNo* lista_criar_lista() {
  ListaNo *q = (ListaNo*) malloc(sizeof(ListaNo));
  q->dado = 0;    // Contador de tamanho
  q->prox = NULL; // Nenhum paciente ainda
  return q; 
} 

void lista_adicionar_final(ListaNo *p, int x) {
  ListaNo *q, *aux;
  q = (ListaNo*) malloc(sizeof(ListaNo));
  q->dado = x;
  q->prox = NULL;
  
  aux = p; // Comeca a partir do no cabeca
  while (aux->prox != NULL) 
    aux = aux->prox; // Caminha ate o ultimo no existente
    
  aux->prox = q; // Encadeia o novo no no fim
  p->dado++;     // Incrementa o contador na cabeca
}
\end{lstlisting}

\begin{conceitobox}{Por que passar \texttt{ListaNo *p} em vez de \texttt{ListaNo **p}?}
Como o nó cabeça foi alocado na criação da lista e seu endereço na memória Heap \textbf{nunca se altera} durante a vida útil do cadastro, a função recebe apenas \texttt{ListaNo *p}. O ponteiro da cabeça permanece fixo; apenas o encadeamento dos campos internos (\texttt{aux->prox}) e o valor do contador (\texttt{p->dado}) são alterados.
\end{conceitobox}

\begin{lstlisting}[language=C, caption={{Remoção Segura por Valor}}]
void lista_remover_valor(ListaNo *p, int v) {
  ListaNo* q = p; // q e o predecessor do no a ser verificado
  while (q->prox != NULL){
    if(q->prox->dado == v) break;
    q = q->prox;
  }
  if(q->prox == NULL) return; // Paciente nao encontrado
  
  ListaNo* tmp = q->prox;     // Guarda ponteiro do no a desalocar
  q->prox = tmp->prox;        // 'Pula' o no a ser removido
  free(tmp);                  // Libera a memoria da Heap
  p->dado--;                  // Atualiza o contador na cabeca
}
\end{lstlisting}

\newpage

% -------------------------------------------------------------
% SUBSEÇÃO 3.2: LISTA CIRCULAR
% -------------------------------------------------------------
\subsection{Lista Circular Simplesmente Encadeada (\texttt{lista\_circular.c/.h})}

\subsubsection{Conceito do Ponteiro para a Cauda (\texttt{*p})}
Diferente de uma lista linear onde o último nó aponta para \texttt{NULL}, na lista circular o último nó aponta de volta para o primeiro:
\begin{equation}
    \text{último}\to\text{prox} = \text{primeiro}
\end{equation}

A implementação adota a técnica mais eficiente da ciência da computação para listas circulares: \textbf{manter um único ponteiro para o último nó} (cauda \texttt{p}):
\begin{itemize}
    \item O último elemento é acessado em $O(1)$: \texttt{p}.
    \item O primeiro elemento (o plantonista da vez) é acessado em $O(1)$: \texttt{p->prox}.
    \item A rotação de plantão (\texttt{avancar\_lista}) é uma simples atribuição em tempo constante $O(1)$: \texttt{*p = (*p)->prox}!
\end{itemize}

\begin{center}
\begin{tikzpicture}[>=Stealth, node distance=1.5cm]
    \node [draw, circle, fill=orange!20, minimum size=1.1cm] (m1) {\small Med 1};
    \node [draw, circle, fill=orange!20, minimum size=1.1cm, right=1.8cm of m1] (m2) {\small Med 2};
    \node [draw, circle, fill=orange!20, minimum size=1.1cm, below=1.2cm of $(m1)!0.5!(m2)$] (m3) {\small Med 3};

    \draw [->, thick] (m1) to[bend left=25] (m2);
    \draw [->, thick] (m2) to[bend left=25] (m3);
    \draw [->, thick] (m3) to[bend left=25] (m1);

    \node [below=0.4cm of m3, color=red!70!black] (ptr) {\textbf{*p (Cauda)}};
    \draw [->, thick, red!70!black] (ptr) -- (m3);
    
    \node [above=0.4cm of m1, color=blue!70!black] (primeiro) {\textbf{p->prox (Plantonista da Vez)}};
    \draw [->, thick, blue!70!black] (primeiro) -- (m1);
\end{tikzpicture}
\end{center}

\subsubsection{Implementação das Funções da Lista Circular}

\begin{lstlisting}[language=C, caption={{Avanço de Plantão e Inserção no Fim em $O(1)$}}]
// Avanca o plantao da escala medica em tempo O(1)
void avancar_lista(CircularNo **p) {
  if (*p != NULL) 
    *p = (*p)->prox; // A cauda vira o proximo no; o inicio vira o seguinte!
}

// Retorna o ID do medico plantonista da vez em O(1)
int acessar_primeiro(CircularNo *p){
  if(p == NULL) return -1; // Lista vazia
  CircularNo *q = p->prox; // p->prox e sempre o primeiro no!
  return q->dado;
}

// Insere novo medico no fim da escala
void adicionar_final(CircularNo **p, int x) {
  CircularNo *q = (CircularNo*) malloc(sizeof(CircularNo));
  q->dado = x;
  if(*p == NULL) { // Lista vazia: no aponta para si proprio
    *p = q;
    q->prox = q;
  } else {
    q->prox = (*p)->prox; // O novo no aponta para o primeiro no
    (*p)->prox = q;        // A antiga cauda aponta para o novo no
    *p = q;                // O ponteiro da cauda passa a ser o novo no!
  }
}
\end{lstlisting}

\begin{destaquebox}{Por que \texttt{CircularNo **p} (Ponteiro Duplo)?}
Quando a lista circular está vazia, o ponteiro na função \texttt{main} vale \texttt{NULL}. Para que a função \texttt{adicionar\_final} altere o valor do ponteiro da \texttt{main}, ela precisa receber o \textbf{endereço daquele ponteiro} (\texttt{\&medicos}). O desreferenciamento \texttt{*p = q} atualiza a variável original na memória da \texttt{main}.
\end{destaquebox}

\newpage

% -------------------------------------------------------------
% SUBSEÇÃO 3.3: FILA FIFO
% -------------------------------------------------------------
\subsection{TAD Fila — FIFO (\texttt{fila.c/.h})}

\subsubsection{Encapsulamento e Composição}
O TAD \texttt{Fila} encapsula a complexidade dos nós encadeados. Seu tipo de dado exposto na interface é:
\begin{lstlisting}[language=C]
typedef struct {
  int tam;          // Quantidade de pacientes na sala de espera
  CircularNo *L;    // Lista circular interna usada como buffer FIFO
} Fila;
\end{lstlisting}

\begin{conceitobox}{Vantagem Arquitetural da Fila sobre Lista Circular}
Ao utilizar a lista circular internamente:
\begin{itemize}
    \item \textbf{Enfileirar (\texttt{fila\_adicionar})}: chama \texttt{adicionar\_final(\&L, x)} em $O(1)$.
    \item \textbf{Espiar Primeiro (\texttt{fila\_topo})}: chama \texttt{acessar\_primeiro(L)} em $O(1)$.
    \item \textbf{Desenfileirar (\texttt{fila\_remover})}: chama \texttt{remover\_inicio(\&L)} em $O(1)$.
\end{itemize}
Garante-se o comportamento estrito \textbf{FIFO} (Primeiro a Entrar, Primeiro a Sair) com todas as operações essenciais executadas em tempo constante $O(1)$ sem desperdício de memória!
\end{conceitobox}

\begin{lstlisting}[language=C, caption={{Operações da Fila FIFO}}]
Fila* fila_criar() {
  Fila *p = (Fila*) malloc(sizeof(Fila));
  p->tam = 0;
  p->L = criar_lista(); // Inicializa lista circular interna
  return p; 
}

void fila_adicionar(Fila **p, int x) {
  adicionar_final(&((*p)->L), x); // Entra no fim da fila
  (*p)->tam++;
}

void fila_remover(Fila **p) {
  if ((*p)->tam == 0) return;
  remover_inicio(&((*p)->L));      // Sai do inicio da fila
  (*p)->tam--;
}

int fila_topo(Fila *p) {
  return acessar_primeiro(p->L);   // Proximo da fila a ser atendido
}
\end{lstlisting}

% -------------------------------------------------------------
% SUBSEÇÃO 3.4: PILHA LIFO
% -------------------------------------------------------------
\subsection{TAD Pilha — LIFO (\texttt{pilha.c/.h})}

\subsubsection{Encapsulamento e Histórico Reversível}
O TAD \texttt{Pilha} armazena o histórico dos atendimentos médicos realizados:
\begin{lstlisting}[language=C]
typedef struct {
  ListaNo *L; // Lista encadeada com cabeca usada como buffer LIFO
} Pilha;
\end{lstlisting}

O topo da pilha coincide com a primeira posição logo após o nó cabeça da lista encadeada (\texttt{L->prox}):
\begin{itemize}
    \item \textbf{Empilhar (\texttt{pilha\_adicionar})}: insere no início da lista via \texttt{lista\_adicionar\_inicio(p->L, x)} em $O(1)$.
    \item \textbf{Topo (\texttt{pilha\_topo})}: consulta o primeiro nó via \texttt{lista\_acessar\_primeiro(p->L)} em $O(1)$.
    \item \textbf{Desempilhar (\texttt{pilha\_remover})}: remove o primeiro nó via \texttt{lista\_remover\_inicio(p->L)} em $O(1)$.
\end{itemize}

O elemento inserido mais recentemente fica no topo. Quando o usuário executa a opção de ``Desfazer Atendimento'', a operação de desempilhar (\textit{pop}) devolve exatamente o último paciente consultado!

\newpage

% -------------------------------------------------------------
% SEÇÃO 4: O PROGRAMA INTEGRADOR (clinica.c)
% -------------------------------------------------------------
\section{O Sistema Integrador (\texttt{clinica.c})}

O arquivo \texttt{clinica.c} atua como o controlador da aplicação. Na função \texttt{main}, as quatro coleções são instanciadas independentemente:

\begin{lstlisting}[language=C, caption={{Instanciação das Quatro Estruturas no main}}]
int main(void) {
  ListaNo *pacientes   = lista_criar_lista(); // 1. Cadastro permanente
  CircularNo *medicos  = criar_lista();       // 2. Escala de plantao
  Fila *espera         = fila_criar();        // 3. Sala de espera (FIFO)
  Pilha *historico     = pilha_criar();       // 4. Historico (LIFO)
  int opcao;
  // ... loop interativo ...
\end{lstlisting}

\subsection{Análise das Operações de Integração}

\subsubsection{Atendimento de Paciente (\texttt{Opção 10})}
Esta é a operação mais rica do sistema, pois coordena simultaneamente três coleções distintas:

\begin{lstlisting}[language=C, caption={{Coordenação entre Fila, Pilha e Lista Circular}}]
case 10:
  if (fila_tamanho(espera) == 0) 
    printf("Sala de espera vazia.\n");
  else if (tamanho_lista(medicos) == 0) 
    printf("Nao ha medico de plantao.\n");
  else {
    int paciente = fila_topo(espera);         // 1. Identifica o proximo da fila
    int medico   = acessar_primeiro(medicos);  // 2. Identifica o plantonista da vez
    
    fila_remover(&espera);                     // 3. Paciente SAI da sala de espera (FIFO)
    pilha_adicionar(&historico, paciente);     // 4. Paciente ENTRA no historico (LIFO)
    
    printf("Paciente %d atendido pelo medico %d.\n", paciente, medico);
    avancar_lista(&medicos);                   // 5. Escala circular avanca para o prox medico!
  }
  break;
\end{lstlisting}

\subsubsection{Desfazer Último Atendimento (\texttt{Opção 11})}
Demonstra a utilidade prática da estrutura de Pilha (LIFO) para histórico reversível:

\begin{lstlisting}[language=C, caption={{Reversão de Atendimento via Desempilhamento}}]
case 11:
  if (pilha_tamanho(historico) == 0) 
    printf("Nao ha atendimento para desfazer.\n");
  else {
    int paciente = pilha_topo(historico);     // 1. Consulta o ultimo atendido
    pilha_remover(&historico);                 // 2. Desempilha do historico
    fila_adicionar(&espera, paciente);         // 3. Devolve paciente para a fila de espera!
    printf("Atendimento do paciente %d desfeito; voltou para a fila.\n", paciente);
  }
  break;
\end{lstlisting}

\subsubsection{Encerramento Seguro e Liberação de Memória (\texttt{Opção 0})}
Ao finalizar a aplicação, todas as quatro coleções são desalocadas recursiva ou iterativamente, garantindo vazamento zero de memória (\textit{Zero Memory Leaks}):

\begin{lstlisting}[language=C, caption={{Desalocação Completa na Saída}}]
  lista_destruir_lista(&pacientes);
  destruir_lista(&medicos);
  fila_destruir(&espera);
  pilha_destruir(&historico);
  return 0;
\end{lstlisting}

\newpage

% -------------------------------------------------------------
% SEÇÃO 5: RASTREAMENTO PASSO A PASSO (TRACE DE EXECUÇÃO)
% -------------------------------------------------------------
\section{Rastreamento Passo a Passo de um Cenário Real}

A tabela a seguir demonstra a evolução das quatro estruturas de dados durante uma sequência de ações típica no terminal:

\begin{table}[h!]
\centering
\small
\renewcommand{\arraystretch}{1.3}
\resizebox{\textwidth}{!}{%
\begin{tabular}{@{}clcccc@{}}
\toprule
\textbf{Passo} & \textbf{Ação do Usuário} & \textbf{Cadastro (\texttt{ListaNo})} & \textbf{Médicos (\texttt{CircularNo})} & \textbf{Fila (\texttt{Fila})} & \textbf{Pilha (\texttt{Pilha})} \\ \midrule
1 & Cadastrar pac. 101 e 102 & $[101, 102]$ & $\emptyset$ & $\emptyset$ & $\emptyset$ \\
2 & Cadastrar méd. 1 e 2 & $[101, 102]$ & $[1 \circlearrowleft 2]$ & $\emptyset$ & $\emptyset$ \\
3 & Colocar 101 na fila & $[101, 102]$ & $[1 \circlearrowleft 2]$ & $[101]$ & $\emptyset$ \\
4 & Colocar 102 na fila & $[101, 102]$ & $[1 \circlearrowleft 2]$ & $[101, 102]$ & $\emptyset$ \\
5 & \textbf{Atender paciente} & $[101, 102]$ & $[2 \circlearrowleft 1]$ (avançou!) & $[102]$ & $[\mathbf{101}]$ (topo) \\
6 & \textbf{Atender paciente} & $[101, 102]$ & $[1 \circlearrowleft 2]$ (avançou!) & $\emptyset$ & $[\mathbf{102}, 101]$ \\
7 & \textbf{Desfazer atendimento} & $[101, 102]$ & $[1 \circlearrowleft 2]$ & $[102]$ & $[\mathbf{101}]$ \\ \bottomrule
\end{tabular}%
}
\caption{Trace de Estados da Memória durante o Fluxo Clínico}
\end{table}

\begin{destaquebox}{Observações do Rastreamento}
\begin{itemize}
    \item \textbf{No passo 5}: O paciente 101 é atendido pelo médico 1 (primeiro da escala). A fila passa a ter apenas o 102 (disciplina FIFO). A escala médica avança e o médico 2 assume o plantão da vez.
    \item \textbf{No passo 6}: O paciente 102 é atendido pelo médico 2. A pilha recebe 102 no topo (ficando $[102, 101]$). A escala circular completa a volta e o médico 1 torna a ser o da vez.
    \item \textbf{No passo 7}: Ao desfazer, o paciente 102 (que estava no topo da pilha por ter sido o último atendido) é retirado da pilha e recolocado na fila de espera. O histórico volta a conter apenas o paciente 101.
\end{itemize}
\end{destaquebox}

% -------------------------------------------------------------
% SEÇÃO 6: CONCLUSÃO
% -------------------------------------------------------------
\section{Conclusão}

O projeto \textbf{Clínica VidaSaudável} cumpre integralmente os requisitos fundamentais da disciplina de Estruturas de Dados:
\begin{enumerate}
    \item \textbf{Domínio Algorítmico}: Aplicação das complexidades $O(1)$ nas operações de inserção/remoção em filas, pilhas e rotação de plantão circular.
    \item \textbf{Engenharia de Software em C}: Alta modularidade, isolamento de TADs via arquivos de cabeçalho e ocultação da representação interna.
    \item \textbf{Robustez Operacional}: Tratamento correto de entradas do usuário, prevenção de operações em coleções vazias e destruição completa dos nós dinâmicos da Heap.
\end{enumerate}

\end{document}
